%% ============================================================
%% preamble.tex — ARES 프로젝트 코드 분석 종합 보고서
%% XeLaTeX 컴파일: xelatex CodeReport.tex
%% ============================================================

%% ----- 한국어 지원 -----
\usepackage{kotex}
\usepackage{fontspec}
\defaultfontfeatures{Mapping=tex-text}

%% ----- 폰트 -----
%% 기본값: macOS 시스템 폰트(Apple SD Gothic Neo).
%% Windows(개발 정본 환경)에서 컴파일할 경우 아래 두 줄을
%%   \setmainhangulfont{Malgun Gothic}  또는  KoPub Batang
%% 등으로 교체할 것.
\setmainfont{Apple SD Gothic Neo}
\setsansfont{Apple SD Gothic Neo}
\setmonofont{Menlo}[Scale=0.86]
\setmainhangulfont{Apple SD Gothic Neo}[Scale=1.0]
\setsanshangulfont{Apple SD Gothic Neo}
\setmonohangulfont{Apple SD Gothic Neo}[Scale=0.86]

%% ----- 이모지(흑백) 폴백 폰트 -----
%% Apple Color Emoji 등 컬러 비트맵 폰트는 XeLaTeX가 렌더하지 못해 깨진다.
%% OFL 라이선스 Noto Emoji(흑백, 윤곽선)를 저장소(fonts/)에 동봉하여 폴백으로 사용.
%% ----- 이모지(흑백) 렌더링 -----
%% Apple Color Emoji 등 컬러 비트맵 폰트는 XeLaTeX가 렌더하지 못해 깨진다.
%% OFL 라이선스 Noto Emoji(흑백 윤곽선)를 fonts/에 동봉해 사용.
%% xetexko(kotex)는 상위 평면(astral) 이모지를 "한자(hanja)" 부류로 라우팅하므로,
%% 한자 폰트를 Noto Emoji로 지정하면 🚗🔊🚀 등이 자동으로 렌더된다.
%% (본 문서에는 한자가 없어 안전. 한자를 추가할 경우 \adhochanjafont 사용 필요.)
%% BMP 영역 기호(⏱ U+23F1 등)는 라틴 부류로 처리되므로 \emj{...}로 감싼다.
\newfontfamily\emojifont{NotoEmoji-Regular.ttf}[Path=./fonts/, Scale=0.9]
\setmainhanjafont{NotoEmoji-Regular.ttf}[Path=./fonts/, Scale=0.9]
\setmonohanjafont{NotoEmoji-Regular.ttf}[Path=./fonts/, Scale=0.82] % 코드블록 내 이모지
\newcommand{\emj}[1]{{\emojifont #1}}

%% ----- 레이아웃 -----
\usepackage[a4paper, margin=25mm, headheight=22pt]{geometry}
\usepackage{setspace}
\setstretch{1.18}
\usepackage{microtype}
% 정렬 여유: 등폭 토큰의 경미한 overfull 완화
\emergencystretch=3em
\hfuzz=3pt

%% ----- 색상 -----
\usepackage[table]{xcolor}
\definecolor{deepblue}{HTML}{1A3B5D}
\definecolor{accentcopper}{HTML}{B5651D}
\definecolor{rulegray}{HTML}{B8C0CC}
\definecolor{codebg}{HTML}{F5F6F8}
\definecolor{codekw}{HTML}{0B5394}
\definecolor{codecmt}{HTML}{6A8759}
\definecolor{codestr}{HTML}{A31515}
\definecolor{tablehead}{HTML}{E8EDF3}

%% ----- 표 -----
\usepackage{booktabs}
\usepackage{longtable}
\usepackage{array}
\usepackage{tabularx}
\usepackage{multirow}
\renewcommand{\arraystretch}{1.25}
\newcolumntype{L}[1]{>{\raggedright\arraybackslash}p{#1}}
\newcolumntype{C}[1]{>{\centering\arraybackslash}p{#1}}
% API 시그니처용: 등폭 작은 글꼴, 공백/콤마에서 줄바꿈
\newcolumntype{S}[1]{>{\ttfamily\footnotesize\raggedright\arraybackslash}p{#1}}

%% ----- 코드 리스팅 -----
\usepackage{listings}
\lstset{
  basicstyle=\ttfamily\small,
  backgroundcolor=\color{codebg},
  keywordstyle=\color{codekw}\bfseries,
  commentstyle=\color{codecmt}\itshape,
  stringstyle=\color{codestr},
  numberstyle=\tiny\color{rulegray},
  frame=single,
  rulecolor=\color{rulegray},
  framesep=6pt,
  xleftmargin=10pt,
  xrightmargin=4pt,
  breaklines=true,
  breakatwhitespace=false,
  showstringspaces=false,
  tabsize=2,
  columns=fullflexible,
  keepspaces=true,
}
\lstdefinestyle{py}{language=Python}
% listings에는 JavaScript 내장 언어가 없어 키워드를 직접 정의
\lstdefinelanguage{JS}{
  morekeywords={var,let,const,function,return,if,else,for,while,await,async,
    new,this,Promise,true,false,null,window},
  sensitive=true,
  morecomment=[l]{//},
  morecomment=[s]{/*}{*/},
  morestring=[b]',
  morestring=[b]",
}
\lstdefinestyle{js}{language=JS}
\lstdefinestyle{json}{language={}}

%% 짧은 인라인 명령/식별자
\newcommand{\code}[1]{{\ttfamily\small #1}}

%% ----- 제목 스타일 -----
\usepackage{titlesec}
\titleformat{\section}
  {\sffamily\Large\bfseries\color{deepblue}}
  {\thesection}{0.6em}{}
\titleformat{\subsection}
  {\sffamily\large\bfseries\color{deepblue}}
  {\thesubsection}{0.5em}{}
\titleformat{\subsubsection}
  {\sffamily\normalsize\bfseries\color{accentcopper}}
  {\thesubsubsection}{0.5em}{}
\titlespacing*{\section}{0pt}{1.4em}{0.7em}

%% ----- 머리글/바닥글 -----
\usepackage{fancyhdr}
\pagestyle{fancy}
\fancyhf{}
\renewcommand{\headrulewidth}{0.4pt}
\renewcommand{\headrule}{\hbox to\headwidth{\color{rulegray}\leaders\hrule height \headrulewidth\hfill}}
\fancyhead[L]{\small\sffamily\color{deepblue}ARES 코드 분석 종합 보고서}
\fancyhead[R]{\small\sffamily\color{rulegray}\leftmark}
\fancyfoot[C]{\small\thepage}

%% ----- 박스 (요약/주의) -----
\usepackage[most]{tcolorbox}
\newtcolorbox{infobox}[1][]{
  colback=codebg, colframe=deepblue, boxrule=0.6pt,
  arc=2pt, left=8pt, right=8pt, top=6pt, bottom=6pt,
  fonttitle=\sffamily\bfseries\color{white},
  coltitle=white, colbacktitle=deepblue, #1
}

%% 학생 안내 박스(설계 의도·분석 질문·개선 실습). 페이지 넘김 허용(breakable).
\definecolor{teachbg}{HTML}{FBF4EC}
\newtcolorbox{teachbox}[1]{
  breakable, enhanced,
  colback=teachbg, colframe=accentcopper, boxrule=0.7pt,
  arc=2pt, left=8pt, right=8pt, top=6pt, bottom=6pt,
  fonttitle=\sffamily\bfseries\color{white},
  coltitle=white, colbacktitle=accentcopper,
  title={\hspace{2pt}주요 분석 요소 — #1}
}
%% 분석 박스 내부 소제목
\newcommand{\teachhead}[1]{\smallskip\noindent{\sffamily\bfseries\color{accentcopper}#1}\par\nopagebreak}

%% ----- 기타 -----
\usepackage{enumitem}
\setlist{itemsep=2pt, parsep=1pt, topsep=3pt}
\usepackage{graphicx}
\usepackage{float}

%% ----- 다이어그램(TikZ) -----
\usepackage{tikz}
\usetikzlibrary{arrows.meta, positioning, fit, calc, shapes.geometric,
  backgrounds, shapes.misc}
\tikzset{
  >={Stealth[length=2.4mm]},
  box/.style={draw=deepblue, fill=codebg, rounded corners=2pt,
    align=center, font=\small\sffamily, inner sep=5pt, minimum height=8mm},
  boxw/.style={box, text width=26mm},
  dom/.style={draw=deepblue, very thick, fill=tablehead, rounded corners=3pt,
    align=center, font=\small\sffamily\bfseries, inner sep=6pt, minimum height=10mm},
  hw/.style={draw=accentcopper, fill=white, rounded corners=2pt,
    align=center, font=\footnotesize\sffamily, inner sep=4pt, minimum height=7mm},
  store/.style={draw=accentcopper, fill=codebg, align=center,
    font=\footnotesize\sffamily, inner sep=4pt, minimum height=7mm,
    cylinder, shape border rotate=90, aspect=0.25},
  flow/.style={->, draw=deepblue, thick},
  back/.style={->, draw=accentcopper, thick, dashed},
  dep/.style={->, draw=deepblue!70, semithick},
  lbl/.style={font=\scriptsize\sffamily\itshape, fill=white, inner sep=1.5pt},
  grp/.style={draw=rulegray, dashed, rounded corners=4pt, inner sep=8pt},
}
% 다이어그램 캡션 박스 (#1=캡션, #2=라벨; end에서 쓰기 위해 매크로로 저장)
% \label은 tikzpicture 밖, \caption 뒤에 두어야 그림 번호를 잡는다.
\newenvironment{diagram}[2]
  {\def\diagcap{#1}\def\diaglab{#2}\begin{figure}[H]\centering\begin{tikzpicture}}
  {\end{tikzpicture}\caption{\diagcap}\label{\diaglab}\end{figure}}
\usepackage[hidelinks]{hyperref}
\hypersetup{
  colorlinks=true,
  linkcolor=deepblue,
  urlcolor=accentcopper,
  citecolor=accentcopper,
}
\usepackage{caption}
\captionsetup{font=small, labelfont={bf,sf}, labelsep=period}
